

Fertility is falling almost everywhere. In the United States, the period
fertility rate fell from about 2.1 children per woman in 2007 to 1.62 in
2023, and the world as a whole is likely below replacement for the first
time.\footnote{The total fertility rate was 2.12 in 2007
\citep{cboDemographicOutlook2026} and 1.621 in 2023
\citep{ostermanBirthsFinal2023}. On the world as a whole, see
\citet{fernandezVillaverdeNorrickTerraIncognita2026}.} The decline has become a central concern
for policymakers: fewer births mean fewer workers to support each retiree
and slower economic growth.\footnote{The Congressional Budget Office projects
that deaths will exceed births from 2030 \citep{cboDemographicOutlook2026}.
The Social Security Trustees lowered their long-run assumption for the total
fertility rate from 1.90 to 1.75 children per woman \citep{ssaTrustees2026}.
Across 42 European countries over 1960--2022, lower fertility is associated
with slower growth of real GDP per capita \citep{cevikLostGenerations2025}.}
Housing is the largest item in household budgets, a third of all
spending,\footnote{Housing, including utilities and furnishings, accounted
for 33.4 percent of average annual expenditures in 2024
\citep{blsConsumerExpenditures2024}.} and it is often at the center of the
conversation on fertility.\footnote{Many governments now tie housing support to children.
Hungary offers families a home loan at a fixed rate of 3\% and writes off
10 million forints of the debt at the birth of each child from the second
\citep{hungaryCSOKPlus}. Poland repays 20,000 zloty of a family home loan at
the birth of a second child and 60,000 zloty at a third
\citep{polandFamilyRepayment}. South Korea offers subsidized home loans of up
to 500 million won to households with a newborn \citep{koreaNewbornLoan}, and
Singapore sets aside up to 40\% of new public flats for couples with children
\citep{singaporeFPPS}. In the United States, the Harris campaign proposed up
to \$25,000 in down-payment assistance for first-time buyers in 2024
\citep{harrisCampaign2024}, and the Trump campaign's vice-presidential
candidate said that young people are ``delaying having families'' because
``they can't afford a place to raise a family'' \citep{vanceNewsNation2024}.
The 2025 budget law raised the child tax credit to \$2,200 and created
accounts with a \$1,000 federal deposit for children born from 2025 to 2028
\citep{publicLaw119-21}. In a fourteen-country survey, housing is the third
most cited reason for having fewer children than desired, named by 19
percent of respondents \citep{unfpaSWP2025}.} The two are closely linked. Children need
space, and new parents want extra rooms. Young households, who are the ones
having children, have little savings, so paying for that space is hardest
exactly when the decision is made. While the causes of the decline are
debated,\footnote{\citet{kearneyPuzzleFallingUS2022} find no economic or
policy factor that accounts for it. \citet{kearneyLevineWhyFertilityLow2025}
argue for ``shifting priorities'' across cohorts, a change in norms around
work, parenting, gender roles and leisure, shaped in part by new
technologies such as the smartphone and social media.
\citet{doepkeEconomicsFertilityNew2023} emphasize how easily parents can
combine careers and family.} housing and its affordability will play a role
in the future of family formation.

This paper asks how the cost of family space shapes fertility as births
decline, and what housing policy can do.

